% !TeX program = xelatex
% 润色版本（resume-coach）：仅调整表述、结构与排版一致性，不新增任何未提供的事实或指标。
% 编译方式：在仓库根目录执行 xelatex -output-directory=temp/.out temp/main_algorithm.tex
\documentclass[10pt]{article}

\usepackage{hyperref}
\usepackage{xcolor}
\usepackage{calc}
\usepackage{graphicx}
\usepackage{tikz}
\usepackage{fontspec}
\usepackage{fontawesome5}
\usepackage{titlesec}
\usepackage{enumitem}
\usepackage{fancybox}
\usepackage{xeCJK}

\hypersetup{hidelinks}
\urlstyle{same}

%%%%%%%%%%%%%%%%%%%%
% 设置
%%%%%%%%%%%%%%%%%%%%

\providecommand{\ResumeItemSep}{0.12em}
\providecommand{\ResumeTopSep}{0.18em}
\providecommand{\ResumeArrayStretch}{1.12}
\providecommand{\ResumeLineSpread}{1.08}
\providecommand{\ResumeSectionBefore}{0.55em}
\providecommand{\ResumeSectionAfter}{0.28em}

\setlength{\parindent}{0pt}
\setlength{\parskip}{0pt}
\pagenumbering{gobble}
\setlist[itemize]{
    leftmargin=1.35em,
    itemsep=\ResumeItemSep,
    topsep=\ResumeTopSep,
    parsep=0pt,
    partopsep=0pt
}
\setlist[enumerate]{leftmargin=*}
\renewcommand{\arraystretch}{\ResumeArrayStretch}
\linespread{\ResumeLineSpread}

\titleformat{\section}
  {\Large\bfseries\raggedright}
  {}{0em}
  {}
  [{\color{secondarycolor}\titlerule}]
\titlespacing*{\section}{0pt}{\ResumeSectionBefore}{\ResumeSectionAfter}

\usepackage[
    a4paper,
    left=1.1cm,
    right=1.1cm,
    top=0.75cm,
    bottom=0.85cm,
    nohead
]{geometry}

% 字体设置
% 说明：Noto Serif SC 没有原生斜体字型。将 CJK 的 italic 字型映射到常规字形
%（与过去“找不到斜体、默认替换为常规字形”的实际渲染一致），可消除
% “Font shape `TU/NotoSerifSC(0)/m/it' undefined” 告警；如需真正斜体观感，
% 可在此行追加 ItalicFeatures={FakeSlant=0.18}。
\setCJKmainfont[
    Path=fonts/,
    Extension=.otf,
    BoldFont=*-Bold,
    ItalicFont={NotoSerifSC},
]{NotoSerifSC}

% 自定义颜色
\definecolor{primarycolor}{RGB}{200, 60, 60}
\definecolor{secondarycolor}{RGB}{200, 68, 28}

\newlength{\iconwidth}
\setlength{\iconwidth}{1.5em}

\newcommand{\sectionicon}[2]{\makebox[\iconwidth][c]{\color{primarycolor}{#1}}\quad #2}

\newcommand{\projectheading}[3]{%
    \noindent\makebox[\textwidth][s]{%
        \makebox[0pt][l]{{\large \textbf{#1}}}%
        \hfill
        \makebox[0pt][c]{{\normalsize \textbf{#2}}}%
        \hfill
        \makebox[0pt][r]{#3}%
    }%
}

% 通用个人信息：修改这里即可同步到所有岗位模板。
\newcommand{\ResumeName}{无绝}
\newcommand{\ResumeBirthDate}{2002.10}
\newcommand{\ResumeHometown}{湖北咸宁}
\newcommand{\ResumeSchool}{重庆邮电大学}
\newcommand{\ResumeEnglishLevel}{CET-6}
\newcommand{\ResumeEmail}{250xxxxxxx@qq.com}
\newcommand{\ResumePhone}{198xxxxxxxx}
\newcommand{\ResumeHomepageLabel}{github.com/Dithob}
\newcommand{\ResumeHomepageUrl}{https://github.com/Dithob}
\newcommand{\ResumePhoto}{images/xjpic.jpg}
\newcommand{\ResumeHeaderImage}{images/foot.png}

% 通用教育背景：修改这里即可同步到所有岗位模板。
\newcommand{\ResumeGraduateSchool}{重庆邮电大学}
\newcommand{\ResumeGraduateMajor}{计算机技术}
\newcommand{\ResumeGraduateDegree}{硕士}
\newcommand{\ResumeGraduateDate}{2024.09--2027.06}
\newcommand{\ResumeGraduateAwards}{特等学业奖学金 $\times$ 1、二等学业奖学金 $\times$ 1}

\newcommand{\ResumeUndergraduateSchool}{重庆邮电大学}
\newcommand{\ResumeUndergraduateMajor}{计算机科学与技术}
\newcommand{\ResumeUndergraduateDegree}{本科}
\newcommand{\ResumeUndergraduateDate}{2020.09--2024.07}
\newcommand{\ResumeUndergraduateAwards}{三等学业奖学金；GPA: 3.28/4.0（专业前 25\%）}

\newcommand{\ResumeEducationBlock}{%
    {\textbf{\ResumeGraduateSchool}}，\ResumeGraduateMajor，\textit{\ResumeGraduateDegree} \hfill \ResumeGraduateDate
    \begin{itemize}
        \item \textbf{荣誉奖项}：\ResumeGraduateAwards
    \end{itemize}

    \vspace{0.2em}
    {\textbf{\ResumeUndergraduateSchool}}，\ResumeUndergraduateMajor，\textit{\ResumeUndergraduateDegree} \hfill \ResumeUndergraduateDate
    \begin{itemize}
        \item \textbf{荣誉奖项}：\ResumeUndergraduateAwards
    \end{itemize}
}


%%%%%%%%%%%%%%%%%%%%
% 文章内容
%%%%%%%%%%%%%%%%%%%%

\newcommand{\ResumeTargetRole}{AI 应用开发}
\providecommand{\ResumeTargetRole}{求职方向待定}

\newcommand{\ResumeContact}{%
    \footnotesize
    \textcolor{white}{%
        \faEnvelope \quad 邮箱：\href{mailto:\ResumeEmail}{\ResumeEmail}
        \hspace{3.2em}
        \faPhone \quad 电话：\ResumePhone
        \hspace{3.2em}
        \faGithub \quad 主页：\href{\ResumeHomepageUrl}{\ResumeHomepageLabel}
    }%
}

\newcommand{\ResumeContactBar}{%
    \begin{tikzpicture}[remember picture, overlay]
        \node[anchor=north, inner sep=0pt](headercontacts) at (current page.north){%
            \includegraphics[width=\paperwidth]{\ResumeHeaderImage}
        };
        \node[anchor=center] at(headercontacts.center){\ResumeContact};
    \end{tikzpicture}
    \vspace{-1.15em}
}

\newcommand{\ResumeProfileSection}{%
    \section[个人信息]{\sectionicon{\faAddressCard}{个人信息}}
    \begin{tabular*}{\textwidth}{@{\extracolsep{\fill}}lll}
        \textbf{姓\qquad 名}：\ResumeName & \textbf{出生年月}：\ResumeBirthDate & \textbf{籍\qquad 贯}：\ResumeHometown \\
        \textbf{学\qquad 校}：\ResumeSchool & \textbf{英语水平}：\ResumeEnglishLevel & \textbf{求职方向}：\ResumeTargetRole
    \end{tabular*}
}

\newcommand{\ResumeEducationSection}{%
    \section[教育背景]{\sectionicon{\faGraduationCap}{教育背景}}
    \ResumeEducationBlock
}

\newcommand{\ResumeProfileAndEducation}{%
    %%%%%%%%%%%%%%%%%%%%
    % 个人信息与教育背景
    %%%%%%%%%%%%%%%%%%%%

    \begin{minipage}[t]{0.79\textwidth}
        \ResumeProfileSection

        \ResumeEducationSection
    \end{minipage}
    \hfill
    \begin{minipage}[t]{0.14\textwidth}
        \vspace{-0.35em}
        \setlength{\fboxsep}{0pt}
        % 双线框 (\doublebox) 会给照片额外加约 4pt 边框宽；图片若占满 \linewidth 会触发
        % “Overfull \hbox (3.99988pt too wide)” 告警，故图片宽度预留 4pt，使含框总宽恰好等于栏宽。
        \doublebox{\includegraphics[width=\dimexpr\linewidth-4pt\relax]{\ResumePhoto}}
    \end{minipage}
}


\begin{document}

    \ResumeContactBar
    \ResumeProfileAndEducation

    %%%%%%%%%%%%%%%%%%%%
    % 专业技能
    %%%%%%%%%%%%%%%%%%%%

    \section[专业技能]{\sectionicon{\faWrench}{专业技能}}
    \begin{itemize}
        \item \textbf{LLM 应用与 Agent 开发}：熟悉 Dify、LangChain 等 LLM 应用开发框架，能围绕业务设计 RAG 知识问答、工具调用、多 Agent 工作流；熟悉向量检索、BM25 召回与 Rerank 精排等 RAG 关键环节；理解 Harness、Memory、MCP、Skill 等 Agent 核心机制，熟练应用 Claude Code、Codex 等 AI 编程工具。
        \item \textbf{模型微调与推理优化}：理解 Transformer、自注意力等底层架构原理；了解 SFT 与 PPO 等训练范式，熟悉 LoRA 等微调方法，具备 Qwen 微调与模型 API 集成经验，熟悉 vLLM 推理部署以及 KV Cache 管理与连续批处理调优。
        \item \textbf{计算机基础与工程能力}：掌握数据结构、计算机网络与数据库原理；掌握 Python，熟悉 MySQL/SQLite 设计优化与 Redis 缓存架构；熟悉 Linux 常用操作，掌握 Docker 容器化部署与 Git 版本控制；了解测试理论，能用 Postman/JMeter 完成接口联调与压测。
    \end{itemize}

    %%%%%%%%%%%%%%%%%%%%
    % 实习经历
    %%%%%%%%%%%%%%%%%%%%

    \section[实习经历]{\sectionicon{\faBriefcase}{实习经历}}

    \projectheading{腾讯云与智慧产业研发公司}{AI 应用开发}{2026.06--2026.09}
    \begin{itemize}
        \item \textbf{核心职责}：参与腾讯地图 UI 自动化测试 Agent 平台建设，负责 Agent 编排设计、Skill 工作流与 MCP 工具落地及业务域用例生成，打通测试业务全流程。
        \item \textbf{多 Agent 交叉检查与 Skill 评测}：设计“多智能体独立生成 $\rightarrow$ 交叉比对合并”的用例生成流水线，产出 85 条分级用例，结合人工复核形成“AI 生成 + 人工质量判断”闭环；为地图数据查询、地图渲染两类 Skill 规划 120 条评测用例，覆盖能力调用、参数边界、异常输入与安全约束，沉淀 Agent/Skill 能力评测方法论。
    \end{itemize}

    % \vspace{0.25em}

    % \projectheading{重庆啄木鸟网络科技有限公司}{算法工程师}{2025.06--2025.09}
    % \begin{itemize}
    %     \item \textbf{核心职责}：参与家电售后场景的大模型应用研发，覆盖需求分析、数据清洗、模型调研与训练、Prompt/RAG 流程优化、接口设计与测试、FastAPI 服务封装及部署联调等环节。
    %     \item \textbf{小啄 GPT 智能回复模型}：针对多轮客服对话中的上下文丢失、槽位遗漏与知识检索不准问题，设计“短期记忆 + 长期摘要 + 结构化状态”机制，用 JSON 维护家电类型、故障现象、预约时间等关键状态；针对生成环节，对 Qwen 模型进行微调，注入家电维修领域知识与标准客服话术，有效抑制模型幻觉，使最终回复的业务采纳率提升约 10\%；结合 Dense/BM25 双路召回、RRF 融合与 Rerank 精排构建维修知识检索链路，使长尾 Query 检索 Hit Rate 提升约 20\%，提升长轮次对话的流程控制与回复可靠性。
    %     % \item \textbf{智能产品识别与意图理解}：
    %     % 参与智能产品识别、用户意图识别等模块迭代，设计工作流从用户咨询中抽取品牌、型号、品类、故障现象等结构化槽位，并判断对话意图类型；
    %     % 通过问题重写规范口语化/残缺问询，结合语义匹配、RAG 知识库检索与动态追问策略完成候选产品召回、槽位补全和低置信度确认；
    %     % 支撑 1000+ 类产品识别和 20+ 类用户意图实时识别，为后续精准问答、工单分发与人工兜底提供结构化依据。
    % \end{itemize}

    %%%%%%%%%%%%%%%%%%%%
    % 项目经历
    %%%%%%%%%%%%%%%%%%%%

    \section[项目经历]{\sectionicon{\faProjectDiagram}{项目经历}}

    \projectheading{腾讯地图 UI 自动化测试 Agent}{参与框架建设}{2026.06--2026.09}
    \begin{itemize}
        \item \textbf{Agent 架构与多智能体编排}：构建 LLM Agent 质量效能体系，按职责分五层编排 13 个 Skill 与 3 个自研 MCP 服务，打通“自然语言需求 $\rightarrow$ 自动化用例 $\rightarrow$ 知识沉淀”闭环；联动 TAPD/Figma 从需求自动生成可执行 case.py，单用例生成约 10--15 分钟（手写需 1--2 小时），覆盖 17 个业务域、600+ 生成代码单元。
        \item \textbf{MCP 工具开发与多模态感知}：自研设备控制与线上诊断 MCP，接入腾讯地图 MCP 解析坐标构造测试数据；融合 UIAutomator、OCR、OmniParser、OpenCV 四层 UI 感知，链式降级解决自绘控件无 resource-id 的定位问题；scrcpy 长连接投屏将截图耗时从约 1s 降至 50ms。
        \item \textbf{知识工程与语义检索}：设计纯文件知识图谱，基于 23 个研发仓库的 UI 索引做同源静态推理，让 Agent 生成前先读懂研发代码；以 AST 三级过滤扫描 1500+ 私有 helper，结合 jieba + BM25 实现公共方法语义检索与经验回流复用，避免规模化生成时重复造轮子。
        \item \textbf{可靠性工程与自愈闭环}：以意图锁定、门禁探索、九维度自审等护栏机制保障生成质量；通过 Lint 门禁与逐步断言保证代码稳定，以 health 回测通过率淘汰 flaky 用例；失败后基于证据包自动再生成补丁，形成 Self-heal 闭环。
    \end{itemize}

    \vspace{0.35em}
    \projectheading{基于 RAG 的医药数据分析问答助手}{项目主要成员}{2025.10--2026.02}
    \begin{itemize}
        \item \textbf{项目目标}：面向药品采购、零售、库存等全链路数据及 800 余张业务表，基于 Dify、LangChain、Milvus 与 MySQL 建设 LLM 数据分析问答助手，将销售人员的自然语言问题转化为可执行 SQL 与可读分析结论。
        \item \textbf{应用架构设计}：负责 LLM 应用核心链路设计，基于 Dify 完成工作流编排与模型接入，落地“问题理解--Schema 检索--SQL 生成--结果解释--人工反馈”业务链路，降低业务人员理解复杂库表结构的门槛。
        \item \textbf{表结构理解与字段召回}：对数据字典做元数据增强，将表名、字段名、字段含义与常见枚举值拼接为富文本后向量化；结合业务规则过滤、问题类型判断、样例对与 CoT，引导模型优先选择相关表和字段，提高 SQL 生成稳定性。
        \item \textbf{评估与自修正闭环}：引入 SQL 预执行与自纠错机制，捕获 MySQL Error 后回传 LLM 重写；构建约 200 条 Golden Dataset，以 RAGAS 指标、执行准确率与 Bad Case 回流评估检索与 SQL 质量，执行成功率由约 70\% 提升至 90\%+。
    \end{itemize}

    % \vspace{0.35em}

    % \projectheading{多模态家电识别与知识库入库系统}{项目主要负责人}{2025.06--2025.09}
    % \begin{itemize}
    %     \item \textbf{项目目标}：面向家电售后知识库建设场景，针对图片/OCR 识别信息不完整、长尾产品文本特征稀疏和重复入库问题，设计“识别--补全--校验--去重--入库”的自动化流程，降低人工录入和知识维护成本。
    %     \item \textbf{识别与字段补全}：使用 Docker/Miniconda 统一部署 OCR/视觉大模型环境，基于 Dify 编排“图片识别--联网搜索--字段抽取--置信度校验”多 Agent 流程；结合 Query Rewrite 和动态 Prompt 模板补全品牌、型号、规格、品类等产品字段。
    %     \item \textbf{服务封装}：接入 vLLM 优化模型推理服务，并使用 FastAPI 封装识别、补全、去重和入库接口，统一对外提供可复用能力，支撑商品手册、维修手册和客服知识库的持续更新。
    %     \item \textbf{业务落地}：构建 1000+ 类家电产品的自动入库工作流；通过向量匹配、规则校验和二次补全的去重策略，缓解版本迭代导致的知识冗余，降低人工录入与重复维护成本，为下游客服问答与检索系统提供高质量数据基础。
    % \end{itemize}

    %%%%%%%%%%%%%%%%%%%%
    % 在校成果
    %%%%%%%%%%%%%%%%%%%%

    \section[在校成果]{\sectionicon{\faAward}{在校成果}}
    \begin{itemize}
        \item \textbf{竞赛}：2025 “互联网+”大学生创新创业大赛银奖；2025 中国国际大学生创新大赛产业赛道校级二等奖；2025 首届重庆市 AI 大模型创新应用大赛省级三等奖
        \item \textbf{论文}：《LHAF: A Lightweight Hierarchical Adaptive Framework for LLM-based Multimodal Emotion Recognition in Conversations》（SCI 一区在投）
    \end{itemize}

\end{document}
